RAG offers an architectural solution to mitigate several of the most fundamental limitations of LLMs, such as hallucination, the lack of up-to-date data, and the lack of domain-specific knowledge \cite{wu2025ragnlp}. To establish a comprehensive understanding of these concepts, this chapter provides the necessary technical background for both technologies. The first section examines the architectural foundations and inherent drawbacks of LLMs\footnote{This section is partly based on the author's bachelor's thesis \cite{asunmaa2025large}.}, while the subsequent section focuses on the core principles and operational mechanics of plain vanilla RAG.

\subsection{Large-Language Models}
\label{chapter:LLM}

\subsubsection{Evolution and Adoption of Language Models}
Natural language processing (NLP) has traditionally focused on enabling computer systems to understand human-generated language and acquire knowledge \cite{RussellNorvig2010}. Early statistical approaches, such as n-gram language models, attempted to resolve the inherent ambiguity of natural language by predicting the probability distribution of word sequences \cite{RussellNorvig2010}. While these traditional models could perform structured tasks like rule-based information extraction or spam filtering, their capability to comprehend semantic nuances and manage complex linguistic tasks remained limited \cite{LLMSurvey}.

Many pre-trained language models (PLMs), such as GPT-2 \cite{radford2019language}, are based on the transformer architecture \cite{Vaswani2017}, which excels at handling long-range dependencies through parallel training. While early PLMs required task-specific fine-tuning to perform downstream applications \cite{bommasani2021opportunities}, modern LLMs scale this architecture to hundreds of billions of parameters. According to \citet{LLMSurvey}, the transition from PLMs to LLMs is primarily characterized by this massive scale in parameters and training data, which gives rise to emergent abilities, enabling the models to solve complex tasks and follow instructions without explicit task-specific retraining.

Presently, LLMs are among the most widely used technologies behind modern generative AI applications \citep{LLMSurvey}. Even though these models have demonstrated strong performance across various tasks, such as coding \citep{Coignion_2024}, and have been widely adopted\footnote{For instance, data aggregated in the report \cite{StanfordAIIndex2026} highlights a rapid acceleration in corporate AI integration. In Europe, corporate adoption in at least one business function was 91\% in 2025 compared to 57\% in 2023.}, they are not perfect. 

\subsubsection{Architectural and Operational Constraints}

Despite their capabilities, LLMs exhibit several major limitations, including biases, hallucinations, and the ``Lost in the Middle'' (LitM) phenomenon \citep{introgenai, liu2024lost}. Bias refers to the fact that the model’s output heavily depends on its training data, meaning that the generated content may include inappropriate language, conspiracy theories, or gender bias \citep{introgenai}. 

Due to the autoregressive nature of modern LLMs, these models are inherently prone to hallucinations, where they produce outputs that appear plausible but are factually incorrect \citep{ji2023survey}. This occurs because the models lack an intrinsic understanding of facts, but instead, they generate responses by predicting subsequent tokens based on statistical patterns learned during training. This relying mechanism draws entirely from the model's parametric knowledge, meaning the information is permanently embedded within the model's weights during the training phase \citep{bommasani2021opportunities}. This static nature, consequently, makes it computationally and environmentally expensive and impractical to frequently retrain or fine-tune entire models when the underlying information required to complete a task changes \cite{StanfordAIIndex2026}.

While the context windows, i.e., LLMs' input sizes, have grown significantly during the last couple of years, allowing models to process up to millions of input tokens \citep{geminiteam2024gemini, GoogleDeepMind2025Gemini3Flash}, these architectures consistently suffer from the LitM phenomenon. This problem refers to a limitation where models effectively utilize information located at the absolute beginning (primacy bias) or the end (recency bias) of the prompt, while failing to retrieve relevant facts buried within the middle sections \citep{liu2024lost}. In addition, utilizing massive contexts severely increases computational latency and operational expenses, as modern cloud LLM APIs employ token-based pricing models where input costs scale linearly with prompt length and often double when exceeding specific context thresholds \citep{GoogleVertexAIPricing}. These inherent limitations effectively prevent the full utilization of large context windows in production-grade enterprise environments when attempting to ingest massive enterprise knowledge bases directly into the prompt.

Another source of errors in the output is that the training data has a cutoff date, so it does not contain the most recent information and may generate incorrect answers to questions about recent events \citep{dhingra2022time}. To mitigate these limitations, modern LLM-based systems are integrated with web search capabilities to retrieve real-time information from the internet, a paradigm used, for example, by the Google Gemini model family \cite{geminiteam2024gemini}. While this integration helps keep responses up to date, it introduces an external dependency on the reliability of internet sources. Furthermore, relying purely on open-web search does not solve the challenge for enterprise environments, where the required knowledge is locked behind internal intranet pages and private corporate databases rather than publicly accessible web pages.

\subsection{Naive Retrieval-Augmented Generation}

To address the limitations detailed in the previous section, Retrieval-Augmented Generation (RAG) offers an architectural solution that decouples the LLM's reasoning capabilities from its static parametrized memory \citep{lewis2020rag}. RAG enriches the LLMs' parametric knowledge with non-parametric data served dynamically from an external knowledge base. The LLM's output is based on the user's initial prompt, augmented with the data retrieved from the external knowledge base, and therefore forces the LLM to ground its response in the provided context rather than its parametrized data.

The retrieval-augmentation process requires both the source knowledge and the user query to be preprocessed. Preprocessing of the knowledge base contains three phases: chunking, vectorizing, and indexing. During this data preparation phase, source documents---such as PDF files---are chunked into smaller, more manageable segments\footnote{Early architectures like \citet{lewis2020rag} utilized fixed-length segments of 100 words, while modern architectures optimize segment lengths \citep{demourajunior2026adaptive}.}. These chunks are converted into numerical vectors using an embedding model so that their semantic meanings can be mathematically calculated, and they are subsequently stored in an indexed vector database \citep{gao2023retrieval}. During runtime inference, the user's initial prompt is vectorized using the same embedding model. The system then calculates the mathematical distance (e.g., cosine similarity) between the prompt vector and the stored vectors, retrieving a pre-defined number of the most relevant chunks---denoted as the top-$k$---from the database to augment the initial prompt \citep{johnson2019billion}.

To formalize this two-phase pipeline, the retrieval-augmentation process can be represented mathematically. In the first phase, a document corpus is split into $N$ smaller text chunks, forming the external knowledge base $D = \{d_1, d_2, \ldots, d_N\}$. These chunks are then mapped into a vector database using an embedding function $E(\cdot)$, which yields the vector representation $\mathbf{v}_i = E(d_i)$ for each chunk $d_i \in D$.

In the second phase, the user's query $q$ is embedded into a vector form using the same embedding function $E(q)$, after which the top $k$ most relevant chunks $R_q$ are fetched based on their cosine similarity score using the retrieval function $\text{Ret}_k$:
\begin{equation}
    \text{Ret}_k(q, D) = R_q.
\end{equation}
Finally, the query $q$ and the retrieved chunks $R_q$ are fed to the generator model $G(\cdot)$ to produce the final answer $a$, denoted as
\begin{equation}
    G(q, R_q) = a.
\end{equation}

The primary advantage of the RAG architecture lies in its reliance on an external knowledge base. When properly prompted, this mechanism significantly mitigates hallucinations by forcing the LLM to ground its output exclusively on the retrieved context \citep{shuster2021retrieval}. Furthermore, maintaining an external data repository is substantially more cost-effective and computationally practical than continuously retraining or fine-tuning the underlying model \citep{gao2023retrieval}, allowing enterprise systems to remain up-to-date in real time. From a data governance perspective, an external knowledge base ensures that sensitive corporate data remains secure and under internal control \citep{ni2025towards}. Lastly, by selectively retrieving only the most relevant documents, RAG compresses the required input context, thereby reducing API transaction costs and inference latency compared to naive long-context processing.

However, these benefits introduce distinct trade-offs. The major drawback of a naive RAG pipeline is the increased end-to-end latency of the system \citep{gao2023retrieval}. This latency overhead stems primarily from sequential network requests: first querying the embedding model and vector database, and subsequently transmitting the expanded prompt to the LLM. While modern vector distance calculations are highly optimized, the network overhead and the increased time-to-first-token (TTFT) caused by processing larger, augmented prompts add noticeable delay \citep{gao2023retrieval}. Consequently, if the RAG pipeline is not properly optimized, the final prompt can grow excessively large, rendering the production system both financially expensive and operationally slow.

